\setlength{\tabcolsep}{.24cm}
\begin{table*}[!t]
\begin{tabular}{l r  l l l}
Method & Acc. & Paradigm 
	%& \todo{WHAT ABOUT DELETING THIS COLUMNApplicability} 
	& Extra Complexity & Realistic Eval.
\\ \hline
Fix-A-Step (ours) 
	& 85.4 % 85.45 verified in google doc, MT_CombinedAugGrad
	& OOD \emph{helpful}
	%& any $\ell^U$
	& none
	& Heart2Heart
\\
TOOR {\tiny \citep{huangTheyAreNot2022}}
	& $*$78.5 % read from Fig 4 of Huang et al TOOR paper
	& OOD \emph{helpful}
	%& any $\ell^U$
	& Separate NN for OOD discrimination
	& none
\\
CL {\tiny{\citep{cascante-bonillaCurriculumLabelingRevisiting2021}}}
	& $*$83.0
	& OOD harmful % Quote from paper: "We conjecture that our self-pacing curriculum is key to this scenario, where the adaptive thresholding scheme could help filter the out-of-distribution unlabeled samples during training"
	%& pseudo-label $\ell^U$
	& Multiple rounds of training, each from scratch
	& none
\\
OpenMatch {\tiny{\citep{saitoOpenMatchOpensetConsistency2021}}}
	& 82.3 % 82.25
	& OOD harmful
	%& \todo{FixMatch only?}
	& Extra one-vs-all OOD detector / class
	& none
\\
DS3L {\tiny \citep{guoSafeDeepSemiSupervised2020}}
	& 74.7 %74.7 from our implementation 
	% in google doc, Row "DS3L + MT (reported)"
	% sanity check: looks like 77.5 in their Fig 6
	& OOD harmful
	%& any $\ell^U$
	& 3x train time due to bilevel optimization
	& none
\\
MTCF {\tiny \citep{yuMultiTaskCurriculumFramework2020}}
	& 77.0 % read from our Fig 1 (after adding cosine LR, 68.0-->77.0)
	& OOD harmful
	%& any $\ell^U$
	& Extra OOD head, curriculum learning
	& none
\\
UASD {\tiny \citep{chenSemiSupervisedLearningClass2020}}
	& $*$78.2 % in google doc, row "UASD Figure 3 left (reported)"
	& OOD harmful
	%& fixed $\ell^U$
	& none
	& none
\\
Safe-Student
	{\tiny{\citep{heSafeStudentSafeDeep2022}}}
	& $^{\dagger}$n/a
	& OOD harmful
	%& \todo{FILL} %their objective eq11 doesnt fit our eq1?
	& 2 NNs (teacher \& student), extra KL loss
	& none
\end{tabular}
\caption{Comparison of related work on open-set/safe SSL.
\emph{Acc} means accuracy on the CIFAR-10 6-animal task (defined in Sec.~\ref{sec:results-cifar}) with 400 labeled examples/class and an open-set unlabeled set (50\% mismatch).
Fix-A-Step uses a FixMatch base model, as does OpenMatch.
%, or MeanTeacher for older methods.
Numbers come from our implementation except if marked \textbf{*} (copied from cited paper) or $^\dagger$ (not assessed in cited paper).
%\emph{Time:} Runtime for training on 6-animal task.
\emph{Paradigm:} how each method treats out-of-distribution (OOD) images in the unlabeled set, broadly either possibly helpful or likely harmful (and thus in need of filtering).
%\emph{Applicability:} whether the method might be possible with different unlabeled losses $\ell^U$.
\emph{Extra Complexity:} additional neural networks, layers, or runtime concerns that exceed a standard SSL deep classifier like MixMatch.
\emph{Realistic Eval.:} evaluation beyond ``artificial'' unlabeled sets from common datasets like CIFAR, ImageNet, etc.
%$\dagger$: \todo{DELETE? OpenMatch has no realistic SSL evaluation, but evaluates OOD detection on fine-grained data (e.g. breeds of dog).}
%\todo{CITE} and Dogs\todo{CITE}.
}%endcaption	
\vspace{-0.1cm}
\label{tab:related_work_comparison}
\end{table*}
